\section{Preliminaries and Methodology}\label{sec:method}
\paragraph{Setting.} Agents $m=1,\dots,M$ are evaluated on tasks $t=1,\dots,T$ with $R$ runs each, producing scores $Y_{mtr}\in[0,1]$ (binary in our benchmark). We regard tasks as an i.i.d.\ sample from a task population and runs as conditionally independent given agent and task. All inference below is about that population; Section~\ref{sec:limits} discusses what this does and does not license.

\subsection{pass@$k$ and pass$^k$}
With $c$ successes in $n$ attempts of one task, $\binom{c}{k}/\binom{n}{k}$ is an unbiased estimator of $p^k$ (all $k$ attempts succeed), and $1-\binom{n-c}{k}/\binom{n}{k}$ is unbiased for $1-(1-p)^k$ \citep{chen2021codex}. We report the benchmark-level means of these per-task estimators; the plug-in $\hat p^{\,k}$ is biased upwards. pass$^k$ is the reliability metric of tau-bench \citep{yao2024taubench}; it decreases in $k$ whereas pass@$k$ increases.

\subsection{Intervals for a benchmark score}\label{sec:intervals}
Let $p_t$ be the success probability of task $t$, with mean $p$ and variance $\tau^2$ across tasks. The variance of the benchmark mean is
\begin{equation}
\mathrm{Var}(\bar Y)=\frac{\tau^2}{T}+\frac{p(1-p)-\tau^2}{TR},
\end{equation}
while a binomial interval on the pooled $TR$ attempts assumes $p(1-p)/(TR)$. Their ratio is the design effect
\begin{equation}
\mathrm{DEFF}=1+(R-1)\rho,\qquad \rho=\tau^2/\{p(1-p)\},
\end{equation}
so a nominal 95\% pooled interval has coverage $2\Phi(1.96/\sqrt{\mathrm{DEFF}})-1$. We implement Wilson \citep{wilson1927}, Clopper--Pearson \citep{clopper1934} and Beta-Binomial credible intervals (valid for a single task or i.i.d.\ attempts), the \emph{task bootstrap} (resample whole tasks; percentile and BCa \citep{efron1987bca,efron1993bootstrap}), a \emph{runs-only} bootstrap (ignores $\tau^2/T$), a \emph{two-stage} bootstrap (tasks, then runs; counts run noise twice), and an overdispersed beta-binomial profile-likelihood interval for the mean of $p_t\sim\mathrm{Beta}(\mu\kappa,(1-\mu)\kappa)$.

\subsection{Variance decomposition}
We fit the crossed random-effects model $Y_{mtr}=\mu+a_m+b_t+(ab)_{mt}+e_{mtr}$ with variances $\sigma^2_m,\sigma^2_t,\sigma^2_{mt},\sigma^2_e$ by the ANOVA method of moments (expected mean squares in Appendix~\ref{app:ems}), in the spirit of generalizability theory \citep{cronbach1972,brennan2001}. For binary scores this is a linear-probability decomposition: exact for variances of means, not latent-scale. Two derived quantities matter in practice: the variance of the difference between two agents,
\begin{equation}
\mathrm{Var}(\bar Y_a-\bar Y_b)=2\Big(\frac{\sigma^2_{mt}}{T}+\frac{\sigma^2_e}{TR}\Big),
\end{equation}
in which the task main effect cancels, and the generalizability coefficient $E\rho^2=\sigma^2_m/(\sigma^2_m+\sigma^2_{mt}/T+\sigma^2_e/(TR))$, the share of observed between-agent variance that is true skill.

\subsection{Paired comparison and multiplicity}
For agents $a,b$ let $d_t=\bar Y_{at}-\bar Y_{bt}$. The task is the unit of analysis. We provide a sign-flip permutation test (exact up to 14 non-zero differences), a cluster bootstrap for the mean difference, exact McNemar \citep{mcnemar1947} on per-task majority outcomes (the natural test when $R=1$), Cohen's $d_z$ and the probability of superiority $P(d>0)+\tfrac12P(d=0)$. Pairwise p-values are adjusted by Holm \citep{holm1979} or Benjamini--Hochberg \citep{benjamini1995}.

\subsection{Ranking uncertainty}
We resample tasks (the same tasks for every agent), recompute scores and ranks (ties broken uniformly at random) and report rank distributions, $P(\text{top-}k)$, the matrix of $P(\bar Y_i>\bar Y_j)$, and the fraction of pairs that cannot be told apart. A split-half Kendall $\tau$ measures whether a fresh task sample would reproduce the ordering.

\subsection{Power and the tasks-versus-repeats trade-off}\label{sec:power}
Let $s_t=2\sigma^2_{mt}$ and $s_e=2\sigma^2_e$. The mean paired difference has $\mathrm{Var}=(s_t+s_e/R)/T$. At a fixed number of attempts $B=TR$ this equals $(Rs_t+s_e)/B$, \emph{increasing} in $R$: if tasks are free, a fixed number of attempts is best spent on more tasks. With cost $c_\text{task}$ per new task and $c_\text{run}$ per run, minimising the variance at budget $T(c_\text{task}+Rc_\text{run})$ gives
\begin{equation}
R^\star=\sqrt{\frac{s_e\,c_\text{task}}{s_t\,c_\text{run}}}.
\end{equation}
Power uses the noncentral $t$ distribution with $T-1$ degrees of freedom; the planner also reports the minimum detectable effect. Repeats remain essential for reliability metrics and for estimating $\sigma^2_e$ itself.

\subsection{Anytime-valid sequential evaluation}
For per-task quantities $X_t\in[0,1]$ (we use $(d_t+1)/2$) with mean $m$, the capital process $K_t(m)=\prod_{i\le t}(1+\lambda_i(X_i-m))$ with predictable bets $\lambda_i$ is a nonnegative martingale under the hypothesis that the mean equals $m$. By Ville's inequality, $\{m:K_t(m)<1/\alpha\}$ is a confidence sequence valid at all times simultaneously \citep{howard2021cs,waudbysmith2024betting}. We hedge both signs and use the predictable plug-in bet. A comparison stops when the sequence excludes a margin band or lies inside it.
